\section{Conclusion}
\label{sec:conclusion}

A black-box REST API tester can only test what its requests reach. When it
follows only the types in the specification, the API rejects most of its
requests at the first check: on pet-clinic, a type-only generator reached
\NPetclinicBaselineCov of the operations, and all \NMotivBaselineRequests of
its requests to add a pet failed while parsing a random date.

We presented \tool, which parses an OpenAPI specification, generates values
from examples, formats, patterns and names, reuses real ids from earlier
responses, and orders operations by their dependencies. It parsed all
\NSpecs specifications we tried. It kept \NConfTool of its generated values
inside their schema, against \NConfBaseline for a type-only generator. It
raised the operations reached from \NPetclinicBaselineCov to
\NPetclinicFullCov on pet-clinic and from \NGestaohospitalBaselineCov to
\NGestaohospitalFullCov on gestao-hospital with the same request budget.
Values and id reuse gave this gain; operation order changed it by less than
one operation.
